\section{Related Work}
\label{sec:related}

\subsection{Color and multiplexed 2D codes}

Color has been added to 2D codes in two ways. Palette codes give each module one of
several colors and carry reference colors in the symbol: Microsoft's
HCCB~\cite{parikh2008hccb}, HCC2D, which adds palette modules to a QR-style
symbol~\cite{grillo2010hcc2d}, HD2DC~\cite{melgar2019hd2dc}, and JAB Code, which stores its
palette redundantly and is standardized as ISO/IEC 23634~\cite{berchtold2020jab,iso23634}.
Recent JAB Code work calibrates the printed colors to the printer~\cite{bugert2024calib}
and improves version detection for large symbols~\cite{bugert2025version}. Multiplexed
codes instead stack binary layers, one per colorant or color
channel~\cite{bulan2011elliptical,blasinski2013pccc}, so each layer is a standard symbol. Andr\'e and Ferreira multiplexed three and six QR codes in CMY and modeled
the recognition rate on the color constellation~\cite{andre2014colour}; ColorQR streams
RGB-multiplexed QR codes from a phone screen and splits each received frame into its
channels~\cite{toh2016colorqr}; El-Sheref et al.\ fold up to 24 layers into RGB
shades~\cite{elsheref2020fmt}; Huang et al.\ design multiplexed codes whose plain layer a
standard reader can decode~\cite{huang2022mc2d,huang2023halftone}. Other recent schemes
separate the layers by photographing the symbol under red and under blue
light~\cite{sanagustin2025chromatic} or assign
payload, metadata and signature to the three channels~\cite{wahsheh2026rgb}. Most of these
works report capacity or small capture tests. Of the multiplexed QR designs we found, HiQ is
the only one evaluated on a large public set of phone photographs,
CUHK-CQRC~\cite{yang2016hiq,yang2018hiq}; public photographs of other color barcodes also
exist~\cite{lecca2024coco}. Sirmen and \"Ust\"unda\u{g} survey color barcodes and their
sources of distortion~\cite{sirmen2023survey}.

\subsection{Color recovery}

Decoding a color code means assigning each sampled module a color despite illumination,
white balance and channel mixing. Palette codes compare samples with the reference colors
printed in the symbol~\cite{grillo2010hcc2d,berchtold2020jab}; Bagherinia and Manduchi
analyze when colors can be separated without a palette~\cite{bagherinia2011theory}; Ong et
al.\ adapt thresholds and color spaces to the lighting~\cite{ong2011mmcc}. For
per-colorant codes, Blasinski et al.\ model cross-channel interference as a linear mixing
of colorant densities and estimate it per image, either from pilot blocks built into the
synchronization markers or by expectation maximization over the
data~\cite{blasinski2013pccc,bulan2012em}. HiQ trains its classifiers offline: QDA and a
layered SVM, each optionally with a cross-module interference (CMI) term, are fitted to
labelled CUHK-CQRC photographs~\cite{yang2018hiq}. Benito-Altamirano et al.\ use the
code the other way round: they embed reference colors in back-compatible QR codes and use
them to correct the colors of the rest of the scene~\cite{benito2023backcompat}.

Our calibration is closest to the pilot-block estimate of Blasinski et al.: both fit the
color model to each photograph from modules whose colors are known. We differ in the model
and in where the labels come from. We fit class-conditional Gaussians with shrunk
covariances~\cite{friedman1989rda}, which absorb illumination and channel mixing without a
colorant model, and we take the labels from the finder, alignment and timing patterns whose colors the
symbol format already fixes, so the method runs on existing designs such as HiQ without
changing them, and no photographs are needed for training.

\subsection{Blur and inter-module interference}

When modules shrink toward the size of the optical and demosaicing blur, each sampled
module mixes with its neighbors. Bagherinia and Manduchi model a blurred patch as a linear
combination of its own color and its neighbors' colors and decode all patches jointly by
iteratively minimizing the observation error~\cite{bagherinia2014blur}. HiQ named the effect
cross-module interference in high-density color QR codes and learns neighbor weights
jointly with its classifier~\cite{yang2018hiq}. For monochrome QR codes, blind deblurring
estimates a blur kernel in the pixel domain: from the finder patterns as a known
signal~\cite{vangennip2015deblur}, from priors on the finder and alignment patterns of each
quadrant~\cite{zheng2024positioning}, or by a multi-scale search that stops once the code
reads~\cite{chen2021scale}. Channel models for print-to-camera and screen-to-camera links
estimate spatial and color distortion from training symbols~\cite{zhang2019printcapture} or
equalize color crosstalk with an MMSE filter on a transform-domain
code~\cite{fujihashi2021mimo}.

Our leak model~\eqref{eq:leak} is a linear neighbor model of the same kind as Bagherinia and
Manduchi's, with four edge neighbors and one shared weight. Fixing that weight gives a
closed-form inverse on the module grid, computed in one pass with no kernel estimate and no
iteration over the symbol.

\subsection{Learned decoders and screen-to-camera codes}

Deep networks are now common for barcode localization and are used for decoding as
well~\cite{wudhikarn2022review}. Liao et al.\ train CNNs on pairs of adjacent modules to read
laser-marked Data Matrix symbols~\cite{liao2022datamatrix}. Most learned color decoders come
from screen-to-camera and optical camera communication, where the transmitter is a display:
Deep D2C-Net and DeepLight decode data embedded in displayed images with trained
networks~\cite{tamang2021deepd2c,tran2021deeplight}, DeepCCB-OCC applies deep learning to
complementary color barcodes~\cite{kim2022deepccb}, and CNN demodulators and neural
equalizers read color shift keying with 4 to 512 colors~\cite{gordillo2022cnncsk,onodera2023csk512}.
Model-based display codes use per-channel encoding~\cite{dinesh2019channelwise} or
temporal color vibration~\cite{abe2020vibration}. These systems are trained or tuned for a
display and camera and often combine several frames. Our setting is one photograph of a
printed or displayed static symbol taken with an unknown phone, and our method uses no
training data. We found no learned decoder for multiplexed color QR codes in this period
that reports results on CUHK-CQRC.
